\subsection{Optimal execution as stochastic control: the benchmark for learned policies}

\begin{qfmodel}[Optimal execution with order flow as stochastic control: the Cartea--Jaimungal objective, dynamics and dynamic-programming equation, as presented by Lehalle and Laruelle]\label{mod:cartea-jaimungal-execution}
\qffield{Statement} A trader sells an inventory $Q_0$ by time $T$, choosing a trading rate $\nu_t$ \citep[app.~A, sec.~A.7.3, pp.~281--282]{alpha-laruelle}. The state is the midprice $S_t$, the net order flow $\mu_t$ of all other traders, the remaining inventory $Q_t$ and the cash account $X_t$, with
\[ dS_t=b\,(\mu_t-\nu_t)\,dt+\sigma\,dW_t,\qquad \hat S_t=S_t-\Bigl(\frac\psi2+\kappa\nu_t\Bigr),\qquad dQ_t=-\nu_t\,dt,\qquad dX_t=\hat S_t\,\nu_t\,dt \]
(p.~282): permanent impact linear in the trading rate, and temporary impact $\kappa\nu_t$ with full immediate decay (p.~281). The reward of a given policy $\nu$ is
\[ H_t(\nu)=\mathbb E_t\Bigl[X_T^\nu+Q_T^\nu\Bigl(S_T-\frac\psi2-A\,Q_T^\nu\Bigr)-\phi\int_t^T(Q_s^\nu)^2\,ds\Bigr]. \]
This is the book's objective (p.~281) with the half-spread added to the terminal payoff, so that it agrees with the book's own terminal condition (Interpretation; the printed objective is discussed under Pitfalls). The value function is the best such reward over admissible policies, $H^*=\sup_\nu H(\nu)$, finite in the parameter regime given under Conditions. By dynamic programming it formally solves (Interpretation: the book's Proposition~6, p.~282, with the corrections listed under Pitfalls)
\[ 0=\Bigl(\partial_t+\tfrac12\sigma^2\partial_{SS}+\mathcal L^\mu\Bigr)H^*-\phi q^2+\sup_\nu\Bigl\{\nu\Bigl(S-\frac\psi2-\kappa\nu\Bigr)\partial_xH^*+b(\mu-\nu)\,\partial_SH^*-\nu\,\partial_qH^*\Bigr\}, \]
with $H^*(T,x,S,\mu,q)=x+q(S-\psi/2)-Aq^2$, the terminal payoff above, where $\mathcal L^\mu$ is the generator of the flow. Defining $h:=H^*-x-q(S-\psi/2)$, so that $h(T,\mu,q)=-Aq^2$, gives
\[ \partial_th+\mathcal L^\mu h+b\mu q-\phi q^2+\sup_\nu\bigl\{-\kappa\nu^2-(bq+\partial_qh)\nu\bigr\}=0,\qquad \nu^*=-\frac{bq+\partial_qh}{2\kappa} \]
(p.~283; the half-spread cancels in this decomposition). With $h=h_0+qh_1+q^2h_2$, the coefficient $h_2$ solves a Riccati equation that does not involve $\mu$, $h_1$ carries the expected future flow, and $h_0$ follows by integration (pp.~283--284). The authors' point is that the trader can act on a view of future flows (p.~284). Interpretation: with $b+2h_2=-2\xi$, the optimal rate is
\[ \nu_t^*=\frac{\xi(t)}{\kappa}\,Q_t-\frac{h_1(t,\mu_t)}{2\kappa},\qquad h_1(t,\mu)=b\int_t^Te^{-\frac1\kappa\int_t^s\xi(u)\,du}\,\mathbb E_{t,\mu}[\mu_s]\,ds, \]
so the trader sells in proportion to inventory and sells more slowly when other traders are expected to keep buying.

\qffield{Notation} Units, for a sale: $Q$ in shares; $\nu$ and $\mu$ in shares per unit time, $\nu>0$ selling and $\mu>0$ net buying by others; $S$, $\psi$ and $\hat S$ in currency per share; $b$ in currency per share$^2$, the price change per share of net flow; $\sigma$ in currency per share per square root of time; $\kappa$ in currency $\times$ time per share$^2$; $\phi$ in currency per share$^2$ per unit time; $A$ in currency per share$^2$. Then $h$ is in currency, $h_1$ in currency per share, $h_2$ and $\xi$ in currency per share$^2$, and $\gamma=\sqrt{\phi/\kappa}$ per unit time; $A-b/2$ and $\sqrt{\kappa\phi}$ are both in currency per share$^2$, as the formula for $\zeta$ requires.

Signs: the trader's own selling lowers $S$ permanently by $b$ per share; each fill is $\psi/2+\kappa\nu$ below the midprice; others' net buying raises $S$. $H_t(\nu)$ is the reward of one fixed policy; the value function is $H^*$.

\qffield{Conditions} Model timeless; the survey's account of RL as of February 2023. The trader observes $S$, $\mu$, $Q$ and $X$ at each time. The flow $\mu$ is Markov with a known generator and does not respond to the trader. The parameters $b$, $\kappa$, $\sigma$, $\psi$, $\phi$ and $A$ are constant: the book writes $b_t$ in the price dynamics but a constant $b$ in the equation.

\emph{Finite value.} Interpretation: with the cash account starting at $X_0=0$, so that cash is measured from the start of execution (otherwise add $X_0$ to the right-hand side), integrating by parts, the realised reward of any policy is
\[ \begin{aligned} &X_T+Q_T\Bigl(S_T-\frac\psi2-AQ_T\Bigr)-\phi\int_0^TQ_t^2dt\\ &\quad=Q_0S_0-\frac{\psi Q_0}2-\frac{bQ_0^2}2+b\int_0^TQ_t\mu_t\,dt+\sigma\int_0^TQ_t\,dW_t-\kappa\int_0^T\nu_t^2dt-\phi\int_0^TQ_t^2dt-\Bigl(A-\frac b2\Bigr)Q_T^2, \end{aligned} \]
because the trader's own permanent impact contributes $b(Q_T^2-Q_0^2)/2$ whatever the path. A simple sufficient regime for a finite value and for the formulas used here is $\kappa>0$, $\phi>0$ and $A>b/2$, with admissible policies progressively measurable with respect to the observed state and $\mathbb E\int_0^T\nu_t^2dt<\infty$, and a flow with $\mathbb E\int_0^T\mu_t^2dt<\infty$. Then $\mathbb E\int Q\,dW=0$; $|\mathbb E\int Q\mu\,dt|\le\|Q\|\,\|\mu\|$ with $\|Q\|\le\sqrt TQ_0+T\|\nu\|$ (norms in $L^2(\Omega\times[0,T])$), so the expected reward is bounded above by a concave quadratic in $\|\nu\|$; and $|\zeta|<1$, so $\xi$ is finite and positive on $[0,T]$. The regime is sufficient, not necessary, and no verification theorem is proved here: the quadratic ansatz is a formal solution. Outside it the value can be infinite.

With $T=Q_0=1$, $\mu=\sigma=\psi=0$, $b=0.1$, $\kappa=A=0.01$, $\phi=0.001$ and a constant buying rate $\nu=-M$, the reward's $M^2$ coefficient is $b/2-\kappa-A-\phi/3\approx0.0297>0$: buying raises the price at which the enlarged position is marked at $T$ faster than the penalties grow. Simulated rewards divided by $M^2$ converge to this coefficient.

\emph{Constraints and completion.} The unconstrained problem is a formal benchmark. Economically valid one-directional execution requires $\nu\ge0$, and $Q_t\ge0$. The unconstrained optimum can ask the seller to buy back, for instance when strong expected buying flow makes $h_1$ large, and the half-spread term then credits the purchase (Pitfalls). Clipping $\nu^*$ at zero does not solve the constrained problem: the constraint changes the value function itself, whose equation then takes the supremum over $\nu\ge0$ and has a different $h$. A finite $A$ charges the inventory left at $T$, which is sold at $S_T-\psi/2-AQ_T$. That is a quadratic penalty, not mandatory completion; completion is forced only in the limit $A\to\infty$.

\emph{Penalty and variance.} The running penalty $\phi\int Q^2$, which the book calls risk aversion (p.~282), is not in general the variance of implementation shortfall. For a deterministic schedule $Q$ the random part of the reward is $b\int Q_t\mu_t\,dt+\sigma\int Q_t\,dW_t$. Its variance is
\[ b^2\operatorname{Var}\Bigl(\int Q_t\mu_t\,dt\Bigr)+\sigma^2\int Q_t^2\,dt+2b\sigma\operatorname{Cov}\Bigl(\int Q_t\mu_t\,dt,\int Q_t\,dW_t\Bigr), \]
and the covariance is nonzero when flow and price shocks are correlated. Thus $\sigma^2\int Q^2dt$ is only the contribution of the Brownian innovations. It is the whole variance when $b=0$ or the flow path is deterministic; only then does $\phi=\lambda\sigma^2$ reproduce a mean--variance penalty. Example: $Q_t=1-t$, $\mu_t=B_t$ a Brownian motion, $T=1$, $\sigma=0$ give variance $b^2\int_0^1\!\int_0^1(1-s)(1-t)\min(s,t)\,ds\,dt=b^2/20$ (exact). For a policy that adapts to the flow, $\phi$ is a penalty on holding inventory, not a variance.

\emph{Benchmark limit.} With $\mu\equiv0$ and $A\to\infty$ the optimal holding is $Q_t=Q_0\sinh(\gamma(T-t))/\sinh(\gamma T)$ with $\gamma=\sqrt{\phi/\kappa}$. This is the continuous-time form of the \hyperref[mod:almgren-chriss-execution]{Almgren--Chriss schedule}, with $\kappa$ in the role of $\eta$ and $\phi$ of $\lambda\sigma^2$ (Interpretation; matched numerically to four decimals).

\qffield{Practice} The parameter intuition holds under the model's assumptions, with qualifications (Interpretation; own numerical paths with $\mu\equiv0$, $Q_0=1$, $T=1$, base case $\kappa=0.01$, $\phi=0.001$, $A=0.5$, in consistent units):
\begin{itemize}
\item a larger $\phi$ front-loads the sale: tenfold $\phi$ leaves 0.45 rather than 0.50 at $T/2$;
\item a larger $A$ mostly cuts what is left at $T$, 0.2 rather than 1.9 percent for tenfold $A$, and barely changes the early path;
\item a larger $\kappa$ slows and smooths trading: tenfold $\kappa$ leaves 0.58 at $T/2$, and with finite $A$ it also leaves 17 percent unsold at $T$.
\end{itemize}
With flow, the rate also depends on the expected $\mu$. Lehalle and Laruelle suggest Almgren--Chriss when price formation is well understood at a 5--30-minute scale, and this framework when one knows a continuous-time version of price formation focused on trading flows (p.~284).

As a benchmark for learned policies: stochastic control solves for a policy given a specified model. Simple models are tractable but may oversimplify the market; realistic ones are often intractable \citep[sec.~1, p.~2]{wide-hambly-xu-yang}. Reinforcement learning learns decisions from repeated interaction with a system, without exemplary supervision or a complete model of the environment (sec.~2, p.~3). When the transition law and rewards are unknown the problem becomes one of learning: \emph{model-based} methods learn them explicitly, \emph{model-free} methods implicitly (sec.~2.2, p.~7). The survey motivates RL for execution by Almgren--Chriss's reliance on its dynamics and linear impact, and by its pre-planned schedule, which ignores real-time conditions (sec.~4.2, p.~33).

The review also reports one study that combined the Almgren--Chriss solution with Q-learning, and one double-DQN study that added a final one-second step to force liquidation (p.~34). On data, most RL algorithms need more samples than the relevant historical data provide, and financial series are non-stationary, hence the call for sample-efficient methods or good market simulators (sec.~5, p.~49). Interpretation: training against a simulator is interaction, but the simulator must be validated. Replaying historical order books does not show how the market would have responded to the agent's own orders, so own impact must still come from a model. Learning from logged historical decisions without interaction is a different setting.

\emph{A second review}. Kolm and Ritter frame execution as one of the intertemporal-choice problems reinforcement learning can solve ``in an almost model-free way''. Their Example~3.2 writes the Almgren--Chriss mean--variance objective as the value function to maximise \citep[undated preprint, Internet Archive capture of 28 May 2024, abstract and Example~3.2, pp.~1 and~9]{wide-kolm-ritter-rl}. They present learning as a replacement for the model-based solution.

\qffield{Pitfalls} Inconsistencies in the appendix (pp.~281--284; corrections are Interpretation):
\begin{itemize}
\item Proposition~6 defines the value function maximising the sale as ``$H^*:=\min_\nu H(\nu)$'' while its equation takes a $\sup$; the value is $\sup_\nu H(\nu)$.
\item The printed objective is $H_t(\nu)=\mathbb E_t[X_T^\nu+Q_T^\nu(S_T-AQ_T^\nu)-\phi\int_t^T(Q_s^\nu)^2ds]$ (p.~281). Its terminal payoff lacks the half-spread that the book's terminal condition $x+q(S-\psi/2)-Aq^2$ (p.~282) charges, and its ansatz and $h$-equations follow the terminal condition. This entry therefore uses the corrected payoff of the Statement. There every share pays $\psi/2$, whether sold before or at $T$: the total charge is the constant $\psi Q_0/2$, and the spread does not affect the control. Under the printed payoff the natural decomposition is a different one, $g:=H-x-qS$ with $g(T,\mu,q)=-Aq^2$, which gives
\[ \partial_tg+\mathcal L^\mu g+b\mu q-\phi q^2+\sup_\nu\Bigl\{-\kappa\nu^2-\Bigl(bq+\partial_qg+\frac\psi2\Bigr)\nu\Bigr\}=0,\qquad \nu^{\mathrm{pr}}=-\frac{bq+\partial_qg+\psi/2}{2\kappa}. \]
Here $g_2=h_2$, and $g_1+\psi/2$ solves the $h_1$ equation with terminal value $\psi/2$, so
\[ \nu^{\mathrm{pr}}_t=\nu^*_t-\frac{\psi}{4\kappa}\,e^{-\frac1\kappa\int_t^T\xi(u)\,du}: \]
shares held to $T$ escape the half-spread, so selling slows, most near $T$. Both policies agree with a direct discretised optimisation with $\mu=\sigma=0$. With $\kappa=0.01$, $\phi=0.001$, $A=0.05$, $b=0.02$, $\psi=0.01$ and $T=Q_0=1$, the inventory left at $T$ is 0.195 under the corrected payoff and 0.245 under the printed one; direct and feedback paths agree to four decimals.
\item The equation prints $-\partial_qH^*$ inside the supremum. The inventory dynamics $dQ=-\nu\,dt$ give $-\nu\,\partial_qH^*$, and the book's own next step on p.~283 is consistent only with $-\nu\,\partial_qH^*$.
\item Notation: p.~282 says ``add the bid--ask spread $\phi$'' where the spread is $\psi$ and $\phi$ is the running penalty. The equation writes $k$ for $\kappa$ and $b_t$ for $b$, and the terminal conditions for $h_0$ and $h_1$ are printed at $t$ rather than $T$.
\item The half-spread enters $\hat S$ linearly in $\nu$, so a negative rate (buying back) would be credited the half-spread. The spread model is economically coherent only for one-directional selling, which the formal benchmark does not impose (Conditions).
\item The printed closed form $h_2=-\xi-b/2$, $\xi(t)=\sqrt{\kappa\phi}\,(1+\zeta e^{2\gamma(T-t)})/(1-\zeta e^{2\gamma(T-t)})$, $\zeta=(A-b/2+\sqrt{\phi\kappa})/(A-b/2-\sqrt{\phi\kappa})$ gives $h_2(T)=A-b$, not $-A$. The Riccati solution has $\xi(t)=\sqrt{\kappa\phi}\,(1+\zeta e^{-2\gamma(T-t)})/(1-\zeta e^{-2\gamma(T-t)})$ with $\zeta=(A-b/2-\sqrt{\kappa\phi})/(A-b/2+\sqrt{\kappa\phi})$; equivalently, the printed expression with the sign of $\xi$ reversed. This agrees with numerical integration.
\item The boxed $h_1$ ``when $b=0$'' omits the factor $b$: with $b=0$ the source term vanishes and $h_1\equiv0$, so the flow drops out. The general solution is the $h_1$ above.
\end{itemize}
This formulation follows Lehalle and Laruelle's presentation of the Cartea--Jaimungal model. Learning a policy or value function from simulated interaction is not the same as learning it from historical data.

\qffield{See also} \hyperref[mod:almgren-chriss-execution]{mean--variance execution} (the discrete deterministic benchmark); \hyperref[def:price-impact]{price impact} (permanent, temporary and transient impact); \hyperref[def:trading-envelope]{trading envelope} (where a schedule sits in a trading algorithm); \hyperref[def:implementation-shortfall]{implementation shortfall}.

\end{qfmodel}
